\documentclass[12pt,a4paper]{article}

\usepackage[utf8]{inputenc}
\usepackage[T1]{fontenc}
\usepackage[czech]{babel}
\usepackage{geometry}
\usepackage{enumitem}
\usepackage{amsmath}
\usepackage{hyperref}
\usepackage{microtype}
\usepackage{graphicx}
\usepackage{float}
\usepackage{tikz}
\usetikzlibrary{patterns}

% --- Spacing (compact, professional) ---
\setlength{\parskip}{3pt}
\setlength{\parindent}{0pt}

\setlist[itemize]{itemsep=1pt, topsep=2pt, parsep=0pt, partopsep=0pt}
\setlist[enumerate]{itemsep=1pt, topsep=2pt, parsep=0pt, partopsep=0pt}

\usepackage{titlesec}
\titlespacing*{\section}{0pt}{6pt}{4pt}
\titlespacing*{\subsection}{0pt}{5pt}{3pt}
\titlespacing*{\subsubsection}{0pt}{4pt}{2pt}

\setlength{\textfloatsep}{8pt}
\setlength{\intextsep}{6pt}
\setlength{\floatsep}{6pt}
\geometry{margin=2.0cm}

\hypersetup{
    colorlinks=true,
    linkcolor=black,
    urlcolor=black
}

\title{
\textbf{\Large Benchmark AI nástrojů pro vývoj průmyslové aplikace}\\
{\large 2D Paletizace (Single-Bin, Python 3.10+, PyQt5)}
}

\author{
Roman Parák, JIC Smart Factory
}

\date{\today
}


\begin{document}

\pagestyle{plain}

\maketitle

%-------------------------------------------------
\section{Účel benchmarku}

Cílem benchmarku je objektivně porovnat vybrané AI nástroje při vývoji realistické průmyslové aplikace.

Každý člen týmu implementuje aplikaci řešící problém \textbf{2D paletizace (2D bin-packing)} na jedné paletě s cílem efektivně využít plochu palety. Aplikace bude obsahovat algoritmické jádro a grafické rozhraní v prostředí PyQt5.

Benchmark proběhne na následujících nástrojích:

\begin{itemize}
    \item ChatGPT -- Codex
    \item Cursor
    \item Claude Code
\end{itemize}

Výsledkem bude funkční aplikace, zdokumentovaný vývojový proces a strukturované hodnocení AI nástroje.

%-------------------------------------------------
\subsection{Ilustrace problému 2D paletizace}

Obrázek~\ref{fig:binpacking_schema} schematicky znázorňuje princip
2D paletizace (single-bin). Obdélník reprezentuje paletu o rozměrech
$W \times H$. Menší obdélníky představují jednotlivé položky.
Položky nesmí být v překryvu a musí respektovat minimální
vzdálenost $g$.

\begin{figure}[H]
\centering
\begin{tikzpicture}[scale=0.8]

% Paleta
\draw[thick] (0,0) rectangle (12,8);
\node at (6,8.5) {Paleta ($W \times H$)};

% Položky
\draw[fill=gray!20] (0.5,0.5) rectangle (4.5,2.5);
\draw[fill=gray!20] (5,0.5) rectangle (8.5,3.5);
\draw[fill=gray!20] (0.5,3) rectangle (3,6);
\draw[fill=gray!20] (3.5,4) rectangle (7,7);

% Otočená položka
\draw[fill=gray!20] (7.5,4) rectangle (10.5,6);

% Neumístěná položka (mimo paletu)
\draw[dashed] (13,3) rectangle (15,5);
\node at (14,2.5) {Neumístěno};

% Označení mezery g
\draw[<->] (4.5,1.5) -- (5,1.5);
\node at (4.75,1.9) {$g$};

\end{tikzpicture}
\caption{Schematická ilustrace problému 2D paletizace (single-bin).}
\label{fig:binpacking_schema}
\end{figure}

%-------------------------------------------------
\section{Funkční požadavky}

\subsection{Účel a rozsah řešení}

Úkolem je implementovat řešení 2D paletizace na jedné paletě (single-bin).

Cílem je:

\begin{itemize}
    \item zabránit překryvům položek,
    \item dodržet minimální mezeru mezi položkami,
    \item zajistit umístění položek uvnitř palety,
    \item maximalizovat využití plochy palety.
\end{itemize}

Řešení může být heuristické. Není požadováno nalezení globálního optima.

%-------------------------------------------------
\subsection{Vstupní parametry}

\textbf{Paleta (Bin)}

\begin{itemize}
    \item $W$ -- šířka palety [mm]
    \item $H$ -- výška palety [mm]
    \item $g$ -- minimální mezera mezi dvěma položkami [mm]
\end{itemize}

Omezení:

\begin{enumerate}
    \item Každá položka musí být umístěna uvnitř oblasti $W \times H$.
    \item Minimální vzdálenost mezi libovolnými dvěma položkami musí být alespoň $g$.
    \item Mezera mezi položkou a hranou palety není vyžadována.
\end{enumerate}

\textbf{Položky (Items)}

Pro typy položek $i = 1, \dots, n$:

\begin{itemize}
    \item $w_i$ -- šířka footprintu [mm]
    \item $h_i$ -- výška footprintu [mm]
    \item $q_i$ -- počet kusů daného typu (libovolné nezáporné celé číslo)
    \item $rotation_i$ -- povolení rotace o 90°
\end{itemize}

Je možné definovat libovolný počet kusů stejného typu položky. Aplikace musí umožňovat také scénář s jediným typem položky ($n = 1$), což odpovídá strukturované paletizaci.

Pokud je rotace povolena, lze použít orientaci $(w_i, h_i)$ nebo $(h_i, w_i)$.

%-------------------------------------------------
\subsubsection*{Validace vstupů}

Aplikace musí:

\begin{itemize}
    \item odmítnout záporné nebo nulové rozměry položek,
    \item odmítnout nekonzistentní vstupy (např. položka větší než paleta
    v obou orientacích),
    \item informovat uživatele o chybě prostřednictvím GUI
    a systémového loggeru,
    \item zabránit spuštění výpočtu při nevalidních datech.
\end{itemize}

%-------------------------------------------------
\subsection{Cíl optimalizace}

Optimalizace je definována lexikograficky:

\begin{enumerate}
    \item maximalizovat počet umístěných kusů,
    \item mezi řešeními se stejným počtem kusů minimalizovat nevyužitou plochu.
\end{enumerate}

Nevyužitá plocha je definována jako:

\begin{equation}
unused\_area = W \cdot H - \sum_i (n_i^{placed} \cdot w_i \cdot h_i)
\label{eq:unused_area}
\end{equation}

Pokud nelze umístit všechny položky, aplikace musí vrátit nejlepší validní řešení a reportovat neumístěné položky (typ a množství).

%-------------------------------------------------
\subsection{Požadovaný výstup}

Aplikace musí poskytovat:

\begin{itemize}
    \item souřadnice každého umístěného kusu,
    \item jeho orientaci,
    \item počet umístěných kusů,
    \item počet neumístěných kusů,
    \item výpočetní čas řešení,
    \item procentuální využití plochy.
\end{itemize}

%-------------------------------------------------
\section{Technické požadavky}

\begin{itemize}
    \item Jazyk: Python 3.10+
    \item GUI framework: PyQt5
    \item Modulární architektura (oddělení GUI a logiky)
    \item Typové anotace a dokumentační řetězce
    \item Veřejný GitHub repozitář
\end{itemize}

%-------------------------------------------------
\subsection{Požadavky na GUI a ovládání aplikace}

Grafické uživatelské rozhraní (GUI) musí být navrženo
přehledně, stabilně a s jasným oddělením vstupní části,
vizualizační oblasti a diagnostického výstupu.

\subsubsection*{1. Zadávání vstupních parametrů}

Aplikace musí umožňovat zadání všech vstupních parametrů
(paleta i položky) přímo prostřednictvím GUI.

GUI musí obsahovat:

\begin{itemize}
    \item vstupní pole pro rozměry palety ($W$, $H$),
    \item vstupní pole pro minimální mezeru $g$,
    \item možnost definovat více typů položek ($w_i$, $h_i$, $q_i$, $rotation_i$),
    \item možnost přidávání a odebírání typů položek,
    \item podporu scénáře s jediným typem položky (strukturovaná paletizace).
\end{itemize}

Alternativně musí aplikace umožňovat načtení vstupních dat
z konfiguračního souboru (např. ve formátu JSON).
Formát konfiguračního souboru musí být jednoznačně popsán
v dokumentaci (README).

Při zadávání parametrů musí být prováděna validace vstupů.
V případě nevalidních dat nesmí být umožněno spuštění výpočtu
a uživatel musí být informován prostřednictvím GUI
a systémového loggeru.

\subsubsection*{2. Ovládací prvky}

GUI musí obsahovat následující základní ovládací prvky:

\begin{itemize}
    \item \textbf{Run / Start} – spustí výpočet paletizace pro aktuální vstupy,
    \item \textbf{Reset Scene} – vymaže aktuální vizualizaci a vrátí aplikaci do výchozího stavu,
    \item \textbf{Clear Logger} – vymaže obsah systémového logu.
\end{itemize}

Ovládací prvky musí být aktivní pouze v odpovídajícím stavu aplikace
(např. nelze spustit výpočet při nevalidních datech).

\subsubsection*{3. Vizualizační oblast}

GUI musí obsahovat grafickou oblast, ve které je zobrazena:

\begin{itemize}
    \item paleta ($W \times H$),
    \item jednotlivé umístěné položky,
    \item jejich orientace,
    \item případné vizuální odlišení neumístěných položek.
\end{itemize}

Vizualizace musí být čitelná, měřítkově konzistentní
a musí reflektovat skutečné výstupy algoritmu.

\subsubsection*{4. System Logger}

Aplikace musí obsahovat systémový log (samostatný panel v GUI),
který zaznamenává běh aplikace a slouží k diagnostice.

Logger musí:

\begin{itemize}
    \item zobrazovat časovou značku a úroveň zprávy (INFO / WARN / ERROR),
    \item zaznamenávat klíčové události (načtení dat, spuštění výpočtu, dokončení, metriky),
    \item zaznamenávat varování (např. částečné řešení, neumístěné položky),
    \item zaznamenávat chyby a výjimky včetně stručného popisu,
    \item podporovat reset obsahu pomocí tlačítka \textbf{Clear Logger}.
\end{itemize}

\subsubsection*{5. Ošetření chyb a stabilita}

GUI musí být robustní a nesmí padat při běžných chybových stavech.
Veškeré chyby musí být zachyceny a korektně reportovány
uživateli prostřednictvím dialogu nebo loggeru.

Aplikace musí po chybě zůstat ve stabilním a znovu použitelném stavu.

%-------------------------------------------------
\section{Struktura projektu}

Repozitář musí mít strukturu:

\begin{itemize}
    \item App/ -- vstupní bod aplikace (UI)
    \item src/ -- aplikační logika a algoritmy
    \item Data/ -- testovací instance a případné výstupy
    \item Example/ -- benchmark skripty
    \item images/ -- obrázky do README
    \item Tests/ -- automatizované testy
\end{itemize}

Součástí repozitáře musí být také pomocné skripty pro zajištění reprodukovatelného prostředí:

\begin{itemize}
    \item \textbf{install.sh / install.bat} – automatizovaný instalační skript.
    
    Skript musí:
    \begin{itemize}
        \item vytvořit nové Conda prostředí,
        \item nainstalovat Python (min. verze 3.10),
        \item nainstalovat všechny požadované knihovny,
        \item umožnit plně funkční spuštění aplikace bez dalších manuálních kroků.
    \end{itemize}

    \item \textbf{verify.py} – validační skript prostředí.
    
    Skript musí:
    \begin{itemize}
        \item ověřit verzi Pythonu,
        \item zkontrolovat dostupnost všech závislostí,
        \item otestovat import hlavních modulů aplikace,
        \item potvrdit, že aplikaci lze korektně spustit.
    \end{itemize}
\end{itemize}

Cílem těchto skriptů je zajistit plně reprodukovatelné prostředí a bezproblémové spuštění projektu na čistém systému.

%-------------------------------------------------
\section{Obsah README.md}

Součástí každého repozitáře musí být profesionálně zpracovaný soubor
\texttt{README.md} v anglickém jazyce. README představuje primární
vstupní dokument projektu a musí umožnit rychlé pochopení účelu,
architektury a způsobu použití aplikace bez nutnosti studia zdrojového kódu.

README musí obsahovat zejména:

\begin{itemize}
    \item jasný a strukturovaný popis projektu a jeho účelu,
    \item stručný kontext problému 2D paletizace,
    \item vizuální ukázku aplikace (screenshot GUI),
    \item systémové požadavky (verze Pythonu, závislosti),
    \item instalační postup (včetně využití skriptů install.sh / install.bat),
    \item návod ke spuštění aplikace,
    \item popis struktury projektu,
    \item přehled implementovaných algoritmů a jejich principu,
    \item známá omezení a případné předpoklady řešení.
\end{itemize}

README musí být psán technicky přesným a profesionálním jazykem
a musí odpovídat kvalitě běžné open-source dokumentace.

%-------------------------------------------------
\section{Požadované výstupy}

Každý účastník benchmarku je odpovědný za dodání následujících výstupů:

\begin{enumerate}
    \item veřejný GitHub repozitář obsahující kompletní zdrojový kód,
    \item plně funkční aplikaci splňující definované požadavky,
    \item README dle stanovené specifikace,
    \item stručnou odbornou reflexi použití AI nástroje (max. 1 strana A4),
    \item prezentaci výsledků v rozsahu 10--15 minut.
\end{enumerate}

Reflexe použití AI nástroje má shrnout zkušenosti z vývoje
a identifikovat přínosy i omezení nástroje v kontextu
průmyslového vývoje softwaru.

%-------------------------------------------------
\section{Metodika hodnocení}

Hodnocení probíhá individuálně na základě odborného posouzení.
Cílem benchmarku není porovnání účastníků mezi sebou,
ale kvalifikované posouzení jednotlivých AI nástrojů.

Při hodnocení se doporučuje zaměřit zejména na:

\begin{itemize}
    \item efektivitu spolupráce s nástrojem během vývoje,
    \item kvalitu a udržitelnost výsledné architektury,
    \item míru manuálních korekcí oproti generovanému kódu,
    \item schopnost nástroje pracovat s kontextem projektu,
    \item celkovou použitelnost nástroje pro průmyslové projekty.
\end{itemize}

%-------------------------------------------------
\section{Metodika benchmarku}

Benchmark je realizován formou individuální implementace
definovaného zadání pomocí různých AI nástrojů.
Každý účastník zodpovídá za kompletní návrh,
implementaci, dokumentaci a prezentaci řešení.

Cílem benchmarku není soutěž mezi účastníky,
ale získání strukturované zkušenosti
s praktickým využitím jednotlivých AI nástrojů.

%-------------------------------------------------
\section{Cíl benchmarku}

Cílem benchmarku je získat kvalifikovaný a strukturovaný přehled
o schopnostech současných AI nástrojů při vývoji průmyslového softwaru.

Benchmark má poskytnout odpovědi na následující otázky:

\begin{itemize}
    \item Do jaké míry jsou AI nástroje schopny samostatně navrhnout
    a implementovat modulární aplikaci střední složitosti?
    \item Jak kvalitní je generovaný kód z hlediska architektury,
    čitelnosti a rozšiřitelnosti?
    \item Jaká je reálná míra manuální korekce a zásahů?
    \item Jak dobře nástroje pracují s kontextem napříč více soubory?
\end{itemize}

Výsledkem benchmarku má být kvalifikované doporučení
pro případné budoucí nasazení AI nástrojů v rámci
vývoje průmyslových projektů.

\end{document}